% ============================================================================
% Section 3 -- the analytical spine. Contribution 1.
%
% STATUS: first full draft, 2026-09-05. This section survives every branch of
% the GH200 gate (paper/PAPER_PLAN_2026-09-05.md section 4) and is deliberately
% written first. Claims C2, C3, C4.
%
% All numbers here: % src: paper/model/p4_*.json @ 2bfc6ca5e
% ============================================================================

\section{When does host-KV offload pay?}
\label{sec:model}

Before describing a policy we establish what there is to gain, because the
answer is narrow and the boundaries are sharp. This section derives four
closed-form results, reports which of them survived contact with hardware, and
uses the survivors to explain a published result that its own authors leave
unexplained.

\subsection{Setup}

Let $\Xlink$ be the achieved interconnect bandwidth between host memory and the
GPU, $\BW$ the achieved HBM bandwidth, $\HBM$ the device memory, $W$ the bytes
of model weights, and $\Ckv = \HBM - W$ the memory left for KV cache. Let $S$
be the context length and $k$ the KV bytes per token.

A decode step reads every weight once and every resident KV byte once. Writing
$\effw$ and $\effkv$ for the fraction of peak bandwidth each achieves --- we
show below that these differ, and that the second is not a constant --- the
step time at batch $B$ is
%
\begin{equation}
\label{eq:step}
\mathrm{ITL}(B) \;=\; \frac{W}{\BW \cdot \effw} \;+\; \frac{B \cdot S \cdot k}{\BW \cdot \effkv}.
\end{equation}
%
The capacity-limited batch is $B_{\mathrm{cap}} = \Ckv / (S k)$: the point at
which KV, not arithmetic, exhausts the device. A \emph{rider} is a request
whose KV lives in host memory and streams over the interconnect during the
step. The number of riders a step can absorb without lengthening it is
$B_{\mathrm{off}} = \mathrm{ITL} \cdot \Xlink / (S k)$.

\subsection{(a) Context length cancels}

The ratio of riders to residents,
%
\begin{equation}
\label{eq:cancel}
\frac{B_{\mathrm{off}}}{B_{\mathrm{cap}}}
  \;=\; \frac{\mathrm{ITL}\cdot \Xlink / (S k)}{\Ckv / (S k)}
  \;=\; \frac{\mathrm{ITL}\cdot \Xlink}{\Ckv},
\end{equation}
%
is independent of $S$. Long context does not make offloading more attractive:
it shrinks both cohorts equally. This is worth stating because the intuition
runs the other way --- long context is what makes KV cache expensive, so it
seems like the regime where moving KV off-device should pay most. It is not.
What a long context buys the rider in freed HBM, it spends again in per-step
transfer.

\paragraph{Validated on hardware.} Equation~\ref{eq:step} predicts that ITL
depends on batch and context only through their product. Across 54
configurations on two models (Llama-3.1-8B, Qwen2.5-32B) on an A100-80, ITL is
linear in total KV tokens $B \cdot S$ with $R^2 = 0.975$, and the residuals are
unbiased across $S \in \{2048, 4096, 8192\}$ (means $+0.09$, $-0.09$, $+0.00$
ms). $B{=}16,S{=}4096$ and $B{=}8,S{=}8192$ both measure $17.17$ ms.
% src: paper/model/p4_llama8b_a100.json, p4_qwen32b_a100.json @ 2bfc6ca5e
We use differential timing --- $\mathrm{ITL} = [T(N_2) - T(N_1)]/(N_2 - N_1)$
with identical prompts in both runs --- so prefill and warmup cancel exactly.

\subsection{(b) The step time is model-dependent}

At capacity a step reads all of device memory once, which invites the
simplification $\mathrm{ITL}_0 \approx \HBM/\BW$: a quantity fixed by the
device alone. It does not hold. Full-pool ITL on A100-80 is $79.3$ ms for
Llama-8B against $55.7$ ms for Qwen2.5-32B,
% src: paper/model/p4_llama8b_a100.json, p4_qwen32b_a100.json @ 2bfc6ca5e
a $1.4\times$ spread across two models on identical hardware. The cause is that
weights and paged KV do not stream at the same efficiency
(\S\ref{sec:model:params}): a model that is mostly weights reads device memory
faster than one that is mostly KV, so $\mathrm{ITL}_0$ depends on the
weight-to-KV ratio and the applicability gate in (d) is model-dependent rather
than a property of the device.

The transfer budget is still $\mathrm{ITL}\cdot\Xlink$, and most of it is still
bought with \emph{weight-read time}. On GH200, Qwen-32B spends roughly
21\,ms of a 34\,ms step reading weights, during which nothing else contends for
the interconnect. The opportunity is time, not spare bandwidth.

\subsection{(c) The gain law}

Adding $B_{\mathrm{off}}$ riders to $B_{\mathrm{cap}}$ residents multiplies the
served batch by
%
\begin{equation}
\label{eq:gain}
\text{gain} \;=\; 1 + \frac{B_{\mathrm{off}}}{B_{\mathrm{cap}}}
      \;\approx\; 1 + \frac{\Xlink}{\BW}\cdot\frac{\HBM}{\Ckv},
\end{equation}
%
substituting $\mathrm{ITL}_0 \approx \HBM/(\BW\,\bar\eta)$ into
Eq.~\ref{eq:cancel}. The structure is worth reading carefully. The bandwidth
ratio $\Xlink/\BW$ is a handicap --- about $0.15$ even on NVLink-C2C --- and it
cannot be improved by anything a serving system does. The amplifier is
$\HBM/\Ckv$: \textbf{the value of offloading tracks how much of HBM the weights
ate}, not model size and not context length.

That single observation is the section's most useful consequence, because it
is counter to how the field selects evaluation models.

Equation~\ref{eq:gain} is a first-order result and we are careful about what it
does and does not include. It counts requests, not tokens per second, and it
ignores two second-order effects: riders lengthen the step slightly (the link
is not infinitely fast), and each rider costs a small, fixed share of device
memory for the rolling window it is streamed into. Table~\ref{tab:gain}
therefore reports \emph{both} the capacity ratio the closed form approximates
and the throughput gain once those two effects are charged, since a reader who
recomputes Eq.~\ref{eq:gain} by hand should land on the first column and not be
surprised.

\begin{table}[t]
\centering
\small
\begin{tabular}{lrrrrrr}
\toprule
Model (GH200-96) & weights & $\Ckv$ & $\HBM/\Ckv$ & ITL$_0$ & capacity & tput \\
\midrule
Llama-8B fp16        & 16.0\,GB & 76.0\,GB & 1.3 & 46.6\,ms & $\times$1.19 & $\times$1.26 \\
Qwen-32B \textbf{fp8}  & 32.0\,GB & 60.0\,GB & 1.6 & 43.3\,ms & $\times$1.27 & $\times$1.31 \\
Qwen-32B \textbf{fp16} & 64.0\,GB & 28.0\,GB & 3.4 & 34.1\,ms & $\times$1.46 & $\times$1.50 \\
Llama-70B fp8        & 72.7\,GB & 19.3\,GB & 5.0 & 33.0\,ms & $\times$1.68 & $\times$1.71 \\
\bottomrule
\end{tabular}
\caption{A second memory tier is an amplifier on $\HBM/\Ckv$, and the table is
monotone in it. \emph{capacity} is $B_{\mathrm{tot}}/B_0$, which
Eq.~\ref{eq:gain} approximates to within 3\% for Qwen-32B fp16 ($\times$1.51
predicted vs $\times$1.46); the residual is the staging buffer's memory cost.
\emph{tput} additionally charges the step lengthening. \textbf{The two Qwen-32B
rows differ only in weight precision}, which isolates $\HBM/\Ckv$ from every
other architectural variable --- see below. The 70B weight figure is the
\emph{measured} fp8 checkpoint (72.7\,GB across 15 shards), not the nominal
70\,GB.}
% src: paper/model/killer_v2.py, analyze(S=8192, GH200, slo=50, stage=0)
\label{tab:gain}
\end{table}

\paragraph{Weight precision is the clean lever on $\HBM/\Ckv$, and we do not
pull it.} Comparing models of different sizes moves the weight footprint, but
also layer count, KV-head count and architecture, so a trend across 8B, 32B and
70B confounds the quantity the law is about with three others. The clean
instrument is a \emph{single} model at two weight precisions, which moves
$\HBM/\Ckv$ and nothing else --- the two Qwen-32B rows above share an
architecture, a tokenizer and a KV geometry and differ only in that fp8 weights
leave 60\,GB for KV where fp16 leave 28\,GB.

We report those rows as projections rather than measurements. Validating the
$\HBM/\Ckv$ axis directly would have cost hours we judged better spent
comparing against the closest prior system, so the law is validated across
\emph{bandwidth} (\S\ref{sec:eval:cost}) while its amplifier rests on the
parameter fit and on agreement with published results
(\S\ref{sec:model:directkv}). \S\ref{sec:eval:model} states this again where
a reviewer will look for it.

The same observation has a practical edge, which \S\ref{sec:discussion} takes
up: quantizing \emph{weights} is not neutral to a memory system, it actively
removes its headroom, while quantizing \emph{KV} scales both cohorts together
and composes.

\paragraph{The model reproduces measurements it was not fitted to.} Because the
weight and KV terms carry different efficiencies, Eq.~\ref{eq:step} predicts a
full-pool step time that depends on the model, and can be checked against the
A100 measurements of \S\ref{sec:model:params} --- which played no part in
constructing it. It predicts 74.4\,ms for Llama-8B and 51.1\,ms for Qwen-32B,
against 79.3 and 55.7 measured: each within 7\%, and the \emph{spread} between
them --- the quantity a single-efficiency model gets structurally wrong,
predicting equality --- reproduced at $1.46\times$ against $1.42\times$
measured.

Several of these parameters are measured with error, so it is worth asking
which ones the projection actually depends on. For Qwen-32B fp16 on GH200-96 at
8K, perturbing the bandwidth efficiencies barely moves it: a $\pm20\%$ error in
$\effkv$ spans $\times1.47$--$\times1.55$, and $\effw$ over 0.70--0.80 spans
$\times1.47$--$\times1.52$. The reason is structural --- the gain depends on a
\emph{ratio}, so a bandwidth error largely cancels between the step-time
numerator and the capacity denominator.

The parameter that does matter is the \textbf{weight footprint}. Moving it from
60 to 68\,GB spans $\times1.45$--$\times1.57$, because weights sit in
$\Ckv = \HBM - W$ and are cancelled by nothing. This is why we use measured
checkpoint sizes rather than nominal parameter counts --- for a quantized model
the two differ by gigabytes and the error propagates undamped into the headline
--- and it is why halving weight precision is the sharpest lever we have on
$\HBM/\Ckv$.
% src: paper/model/killer_v2.py sensitivity sweep, 2026-09-06

Note the direction of the $\effkv$ term, because it is our main downside risk:
a \emph{more} efficient paged-KV read shortens the step, shrinks the transfer
budget, and \emph{lowers} the gain. The curve above is measured on FA-2-era
kernels, and \S\ref{sec:eval} re-measures it on Hopper for exactly this reason.

\subsection{(d) The applicability gate}

Capacity binds only when the ITL SLO exceeds the full-pool ITL:
%
\begin{equation}
\label{eq:gate}
\mathrm{SLO} \;>\; \frac{W}{\BW\,\effw} + \frac{\Ckv}{\BW\,\effkv}.
\end{equation}
%
Below the gate the scheduler is already refusing requests that HBM could hold,
because admitting them would break latency --- the system is latency-limited,
not capacity-limited, and offloading is then strictly harmful. Following (b),
this gate is model-dependent rather than a property of the device alone.

The gate is a scoping statement, and we state it as a finding rather than
apologise for it: this is a feature for \emph{throughput-oriented} serving ---
agentic pipelines, batch inference, long-form and reasoning workloads --- and
not for interactive chat.

\subsection{Measured parameters}
\label{sec:model:params}

A100-80 PCIe; two models; 41 single-shot configurations plus a
13-configuration $n{=}3$ repeat sweep; distinct random contexts; prefix caching
disabled.
% src: paper/model/p4_llama8b_a100_rep3.json @ 2bfc6ca5e

\paragraph{Weight read is efficient and stable.} $\effw = 0.75$ of peak
($\approx$1450\,GB/s), consistent across both models, all contexts, and robust
to the $\effkv$ assumed at $B{=}1$ (0.739--0.773 over that whole range).

\paragraph{Paged KV read is neither.} $\effkv$ degrades by $1.6\times$ with the
KV working set:

\begin{table}[t]
\centering
\small
\begin{tabular}{lrrrrr}
\toprule
KV working set & 4.3\,GB & 8.6\,GB & 17--34\,GB & 43\,GB & 51.5\,GB \\
\midrule
$\effkv$ & 0.77 & 0.72 & 0.54--0.56 & 0.49 & 0.46 \\
\bottomrule
\end{tabular}
\caption{Paged KV read efficiency is not a constant. Error bars from $n{=}3$
repeats are $\pm0.02$--$0.04$; the degradation is far outside them.}
\label{tab:effkv}
\end{table}

This is the real correction to the model, and it matters beyond this paper:
\textbf{every fixed-bandwidth roofline in this literature mis-predicts the
capacity-bound operating point}, because it is exactly at capacity that
$\effkv$ is worst.

\paragraph{A caveat at kernel granularity.} The claim that only total KV bytes
matter holds at three of four same-byte test points but breaks at 17.2\,GB
($B{=}32,S{=}4096$ gives 0.561 against $B{=}16,S{=}8192$ at 0.747). The cause
is a reproducible anomaly at $B{=}32,S{=}4096$ that also appeared in an
independent earlier run --- an attention-kernel scheduling artifact, most
likely FlashDecoding's split-KV heuristics, not measurement noise ($\pm0.61$\,ms
against a $\sim$3\,ms deviation). The aggregate law holds; the
per-configuration surface has kernel-regime steps.

\paragraph{Compute is not the binding term.} Decode-attention arithmetic
intensity equals the GQA group size --- 8 for 70B, 4 for 8B --- against an
H100 balance point near 116 FLOP/byte, so attention sits 14--29$\times$ below
the compute roof. At batch 8, 70B spends 7\% of its step on arithmetic.
% src: paper/model/compute_check.py @ 2bfc6ca5e
Llama-8B at 4K does cross into compute-bound at batch 142, which is a further
reason it is the wrong model to evaluate on.

\subsection{The model explains a published result}
\label{sec:model:directkv}

DirectKV~\cite{directkv} reports that on Llama-3.1-8B its zero-copy host-KV
attention is \emph{slightly slower} than the HBM-resident SGLang baseline, and
that it wins only on OPT-13B and OPT-30B, where SGLang runs out of memory. The
paper does not explain the inversion.

Equation~\ref{eq:gain} predicts it. Llama-8B leaves $\Ckv \approx 76$\,GB of an
H100/GH200 --- $\HBM/\Ckv \approx 1.26$, so there is essentially nothing to
buy, and any overhead the mechanism carries shows up as a net loss. OPT-30B's
60\,GB of weights leave $\Ckv \approx 32$\,GB, more than doubling the
amplifier. \textbf{The applicability gate and the gain law together account for
which of a published system's configurations win and which lose}, which is the
strongest available evidence that the model captures something real rather
than fitting our own measurements.

\ifdraft
\needsgh{$\effkv$ under FA-3 is the top downside risk and sits in the capacity
denominator: if Hopper's kernels close the paged-KV gap, the step shortens, the
rider budget shrinks, and Table~\ref{tab:gain} falls. Re-measured in
\S\ref{sec:eval} (GH200 plan 4.2), which runs early for exactly this reason.}
\fi
